% ===================== PREAMBULE COMMUN (à copier VERBATIM en tête de CHAQUE .tex) =====================
\documentclass[11pt,a4paper]{article}
\usepackage[utf8]{inputenc}
\usepackage[T1]{fontenc}
\usepackage{lmodern}
\usepackage[french]{babel}
\usepackage{amsmath,amssymb,mathtools}
\usepackage{icomma}
\usepackage{xcolor}
\usepackage{colortbl}
\usepackage{tikz}
\usepackage{pgfplots}
\usepackage[most]{tcolorbox}
\usepackage{enumitem}
\usepackage{array}
\usepackage{booktabs}
\usepackage{multicol}
\usepackage{fancyhdr}
\usepackage{caption}
\captionsetup{labelformat=empty,justification=centering,font=small}
\usepackage[hidelinks]{hyperref}
\usetikzlibrary{arrows.meta,positioning,calc,angles,quotes,patterns,backgrounds,shapes.geometric}
\pgfplotsset{compat=1.18}
\usepackage[a4paper,margin=2.2cm,headheight=26pt,headsep=10pt]{geometry}

\definecolor{primary}{RGB}{219,54,77}
\definecolor{primarydark}{RGB}{150,28,46}
\definecolor{primarylight}{RGB}{250,224,228}
\definecolor{defcol}{RGB}{0,114,178}
\definecolor{thmcol}{RGB}{0,140,120}
\definecolor{excol}{RGB}{0,135,90}
\definecolor{methcol}{RGB}{90,90,100}
\definecolor{courbe}{RGB}{40,90,200}
\definecolor{courbeB}{RGB}{210,60,40}
\definecolor{courbeC}{RGB}{30,150,80}
\definecolor{axiscol}{RGB}{60,60,60}
\definecolor{gridcol}{RGB}{200,205,215}

\tcbset{boxbase/.style={enhanced,breakable,boxrule=0.9pt,arc=2.5pt,
   left=3mm,right=3mm,top=2.3mm,bottom=2.3mm,
   fonttitle=\bfseries,coltitle=white,toptitle=0.6mm,bottomtitle=0.6mm}}
\newtcolorbox{methode}[1][]{boxbase,colback=methcol!7,colframe=methcol,sharp corners,
  title={\textsc{Méthode}\ifx\relax#1\relax\relax\else\textmd{ --- \itshape #1}\fi}}
\newtcolorbox{exemple}[1][]{boxbase,colback=excol!5,colframe=excol,
  title={\textsc{Exemple}\ifx\relax#1\relax\relax\else\textmd{ : \itshape #1}\fi}}
\newtcolorbox{idee}[1][]{boxbase,colback=defcol!6,colframe=defcol,
  title={\textsc{L'essentiel}\ifx\relax#1\relax\relax\else\textmd{ : \itshape #1}\fi}}
\newtcolorbox{piege}[1][]{boxbase,colback=primary!7,colframe=primary,
  title={\textsc{Piège}\ifx\relax#1\relax\relax\else\textmd{ : \itshape #1}\fi}}

\pgfplotsset{repere/.style={axis lines=middle,
    axis line style={-{Stealth[length=2.2mm]},axiscol,line width=0.8pt},
    label style={font=\footnotesize},tick label style={font=\scriptsize},
    every axis x label/.style={at={(current axis.right of origin)},anchor=north west},
    every axis y label/.style={at={(current axis.above origin)},anchor=south east},
    xlabel={$x$},ylabel={$y$},grid=both,minor tick num=0,
    major grid style={gridcol,line width=0.4pt},samples=200,clip=false}}

\pagestyle{fancy}\fancyhf{}
\fancyhead[L]{\small\textcolor{primarydark}{\textbf{Fiche 4} — Les fonctions}}
\fancyhead[R]{\small\textcolor{primarydark}{\textsc{Carnet d'entraînement}}}
\fancyfoot[C]{\small\thepage}
\renewcommand{\headrulewidth}{0.8pt}
\renewcommand{\headrule}{\hbox to\headwidth{\color{primary}\leaders\hrule height \headrulewidth\hfill}}
\usepackage{titlesec}
\titleformat{\section}{\Large\bfseries\color{primarydark}}{\thesection}{0.6em}{}
\titleformat{\subsection}{\large\bfseries\color{primary}}{\thesubsection}{0.6em}{}

\newcommand{\R}{\mathbb{R}}\newcommand{\Z}{\mathbb{Z}}\newcommand{\N}{\mathbb{N}}
\newcommand{\Df}{D_f}
\newcommand{\croit}{\textcolor{thmcol}{\Large$\nearrow$}}
\newcommand{\decroit}{\textcolor{thmcol}{\Large$\searrow$}}

\newcommand{\exercice}[1]{\medskip\noindent{\color{primarydark}\bfseries\large Exercice #1}\par\nopagebreak\smallskip}
\newcommand{\titrecarnet}[3]{%
\begin{center}\begin{tikzpicture}
\node[rectangle,rounded corners=7pt,fill=primary,text=white,minimum width=15cm,
  minimum height=1.7cm,align=center,inner sep=8pt]
  {{\LARGE\bfseries #1}\\[2pt]{\normalsize #2}};
\end{tikzpicture}\end{center}
\begin{center}\itshape\color{primarydark}#3\end{center}\medskip}

% ---- RÈGLES D'USAGE DES MACROS ----
% * \croit et \decroit s'emploient en mode TEXTE (dans une cellule de tableau),
%   JAMAIS entre $...$ (ils contiennent déjà un $\nearrow$).
% * \titrecarnet{Titre}{Sous-titre}{Ligne en italique} : bandeau de titre.
% * \exercice{1} : titre d'exercice.
% * Boîtes : \begin{idee}...\end{idee}, \begin{exemple}[titre]...\end{exemple},
%   \begin{methode}[titre]...\end{methode}, \begin{piege}[titre]...\end{piege}.
% Le corps du document commence après \begin{document}.
\begin{document}
\titrecarnet{Thème 3 — Géométrie analytique des droites}{Tutoriel}{Pente, équation, parallèles/perpendiculaires, distances, milieu, intersection.}

\section*{1. Pente et équation d'une droite}

Une droite non verticale a pour équation $y=ax+b$ : $a$ est la \textbf{pente} (coefficient directeur), $b$ l'\textbf{ordonnée à l'origine} (point $(0;b)$).

\begin{idee}[à partir de deux points $A(x_A;y_A)$ et $B(x_B;y_B)$]
\[ a=\frac{y_B-y_A}{x_B-x_A}, \qquad b=y_A-a\,x_A \ \ (\text{ou } b=y_B-a\,x_B). \]
La pente est un \emph{taux d'accroissement} : de combien monte $y$ quand $x$ augmente de $1$.
\end{idee}

\begin{center}
\begin{tikzpicture}
\begin{axis}[repere,width=8cm,height=4.6cm,xmin=-0.6,xmax=4.4,ymin=-0.6,ymax=4.4,
   xtick={1,2,3,4},ytick={1,2,3,4},samples=2]
  \addplot[courbe,line width=1.2pt,domain=-0.3:4.2]{0.75*x+0.5};
  \filldraw[primary](axis cs:1,1.25)circle(1.6pt)node[anchor=south east,font=\scriptsize]{$A$};
  \filldraw[primary](axis cs:3,2.75)circle(1.6pt)node[anchor=north west,font=\scriptsize]{$B$};
  \draw[dashed,gray](axis cs:1,1.25)--(axis cs:3,1.25)--(axis cs:3,2.75);
  \node[font=\scriptsize,anchor=north]at(axis cs:2,1.25){$x_B-x_A$};
  \node[font=\scriptsize,anchor=west]at(axis cs:3,2){$y_B-y_A$};
\end{axis}
\end{tikzpicture}
\end{center}

\section*{2. Parallèles et perpendiculaires}

\begin{idee}
Soient $d:y=ax+b$ et $d':y=a'x+b'$.
\begin{itemize}[leftmargin=1.5em,topsep=2pt]
  \item \textbf{Parallèles} $\Leftrightarrow$ même pente : $a=a'$.
  \item \textbf{Perpendiculaires} $\Leftrightarrow$ $a\cdot a'=-1$, c'est-à-dire $a'=-\dfrac1a$ (opposé de l'inverse).
\end{itemize}
\end{idee}

\section*{3. Distance, milieu, distance point–droite}

\begin{idee}[dans un repère orthonormé]
\[ \overline{AB}=\sqrt{(x_B-x_A)^2+(y_B-y_A)^2}, \qquad M\Big(\tfrac{x_A+x_B}{2};\tfrac{y_A+y_B}{2}\Big). \]
\[ \operatorname{dist}(A,d)=\frac{\lvert a\,x_A+b\,y_A+c\rvert}{\sqrt{a^2+b^2}}\quad\text{avec } d:\ ax+by+c=0\ \text{(forme \emph{canonique})}. \]
\end{idee}

\begin{piege}
Pour la distance point–droite, la droite doit être en forme \textbf{canonique} $ax+by+c=0$. Ici $a,b$ ne sont \emph{pas} la pente et l'ordonnée à l'origine ! Ex. $y=-2x+3$ devient $2x+y-3=0$, donc $a=2,\ b=1,\ c=-3$.
\end{piege}

\section*{4. Intersection de deux droites}

Deux droites sécantes ($a\neq a'$) se coupent en un point : on résout le système
$\begin{cases} y=ax+b\\ y=a'x+b'\end{cases}$ (par substitution ou addition). La solution $(x;y)$ est le point commun.

\begin{exemple}
$y=2x$ et $y=-2x+3$ : en additionnant, $2y=3\Rightarrow y=\tfrac32$, puis $x=\tfrac34$. Point : $\big(\tfrac34;\tfrac32\big)$.
\end{exemple}
\end{document}
